% Propietario conservado: /Users/ruben/Documents/New project/output/REVISION_ARGUMENTAL_APERTURAS_CIERRES_20260915/VI/delivery/MANUSCRIPT_VI_ES_EN_COMPLETE/source_es/sections/13_fock_pauli_composicion.tex
\section{Graded completion of channels and exclusion principle}
\label{vi:vi:sec:fock-pauli}

The construction of a particle provides a space of channels, their
labels and the operators transporting orientation and memory. The
composition of several particles adds a distinct question: how those
channels are exchanged and how their occupations are counted. Here that
layer is constructed from the data already produced by
APP--TRIT--TPK. Mass and the Hamiltonian can then act on
the resulting space; their values do not enter the proof of
the occupation rules.

The parity used is that of sheet transport in the representation
of a particle. It is not the sign of the charge, the parity of a visible
integer or the third coordinate of TRIT considered in isolation.
The different readings preserve their maps.
Section~\ref{vi:vi:sec:conjugacion} constructs the oriented return and the
spinorial representation; this section develops the multiparticle
algebra corresponding to a declared grading and exchange
rule.

\subsection{Finite channel space and parity character}
\label{vi:vi:subsec:canales-graduados}

Given a depth \(K\), let \(S_K\) be the finite set of
channels of the realization under consideration. Its elements preserve the
route and representation data that still distinguish states:
sheet, orientation, boundary, memory, flavor, color and spin
component when they belong to the sector. The quotient defining \(S_K\)
is taken before constructing occupation; it does not identify two channels
solely through equality of their masses.

The one-particle space is
\begin{equation}
 \mathcal H_K=\bigoplus_{s\in S_K}V_s,
 \qquad\langle e_{s,a},e_{t,b}\rangle=\delta_{st}\delta_{ab},
 \label{vi:vi:eq:espacio-canales-fock}
\end{equation}
where \(V_s\) carries the declared representation multiplicity.
Each \(e_{s,a}\) is a complete mode, not an unlabeled cell of the
atlas. Its positive inner product belongs to that finite
realization and makes explicit the normalization used below.

The parity datum is a compositional character of transport,
\begin{equation}
 p:\operatorname{Mor}(\mathsf P_{\mathrm{TPK}})
 \longrightarrow\mathbb Z/2\mathbb Z,
 \qquad p(\gamma\eta)=p(\gamma)+p(\eta).
 \label{vi:vi:eq:caracter-paridad-fock}
\end{equation}
On the channels in which that grading is realized, the
holonomy involution \(J_p\) satisfies
\(J_pe=(-1)^{p(e)}e\). Therefore,
\[
 \mathcal H_K=\mathcal H_{K,\bar0}\oplus\mathcal H_{K,\bar1},
 \qquad \Pi_{\bar0}=\frac{I+J_p}{2},\quad
 \Pi_{\bar1}=\frac{I-J_p}{2}.
\]
The two projectors are orthogonal because \(J_p^*=J_p\) and
\(J_p^2=I\). The equality \eqref{vi:vi:eq:caracter-paridad-fock} ensures
that composing transports preserves the grading. When a realization
has not supplied that character, the following completion is a
construction relative to its explicit choice, not a statistical
attribution obtained from the appearance of a route.

\subsection{Graded exchange and permutation projectors}
\label{vi:vi:subsec:intercambio-graduado}

For homogeneous vectors, the exchange is defined
\begin{equation}
 \tau_p(v\otimes w)=(-1)^{p(v)p(w)}w\otimes v.
 \label{vi:vi:eq:intercambio-graduado}
\end{equation}
The sign is negative exactly when both modes belong to the
odd sector. Exchange does not modify a mode's labels:
it transports the complete mode from one tensor position to another.

\Needspace{10\baselineskip}
\begin{proposition}[Graded representation of permutations]
\label{vi:vi:prop:permutaciones-graduadas}
The adjacent exchanges \(\tau_{p,j}\) on
\(\mathcal H_K^{\otimes n}\) are unitary and satisfy
\[
\begin{gathered}
 \tau_{p,j}^2=I,\qquad
 \tau_{p,j}\tau_{p,k}=\tau_{p,k}\tau_{p,j}\quad(|j-k|>1),\\
 \tau_{p,j}\tau_{p,j+1}\tau_{p,j}
 =\tau_{p,j+1}\tau_{p,j}\tau_{p,j+1}.
\end{gathered}
\]
They therefore define a unitary representation \(U_p\) of
\(\mathfrak S_n\).
\end{proposition}
\begin{proof}
Applying an exchange twice gives the sign
\((-1)^{2p(v)p(w)}=1\) and restores both vectors. Operators
acting on disjoint pairs commute. On a homogeneous triple
of degrees \(a,b,c\), both sides of the three-exchange
relation produce the same final order and the sign
\((-1)^{ab+ac+bc}\). Each operator permutes an orthonormal basis
multiplying it by signs of modulus one, and is therefore unitary.
\end{proof}

The average
\[
 P_n^p=\frac1{n!}\sum_{\pi\in\mathfrak S_n}U_p(\pi)
\]
is the orthogonal projector onto tensors invariant under graded
exchange. In an even homogeneous space it coincides with the
ordinary symmetrizer; in an odd homogeneous one it coincides with the
antisymmetrizer. If \(U\) denotes the unsigned permutation,
\begin{equation}
 S_n=\frac1{n!}\sum_\pi U(\pi),\qquad
 A_n=\frac1{n!}\sum_\pi\operatorname{sgn}(\pi)U(\pi).
 \label{vi:vi:eq:simetrizadores-fock}
\end{equation}

\begin{proposition}[Exact symmetrization and antisymmetrization]
\label{vi:vi:prop:proyectores-fock}
We have \(S_n^*=S_n^2=S_n\) and \(A_n^*=A_n^2=A_n\).
For \(n\ge2\), \(S_nA_n=0\). Their images are, respectively,
\(\operatorname{Sym}^n\mathcal H\) and \(\bigwedge^n\mathcal H\).
\end{proposition}
\begin{proof}
Taking the adjoint replaces each permutation by its inverse and leaves
the sums invariant. In the square of each average, every
permutation \(\rho\) appears \(n!\) times; in the second average
its coefficient is \(\operatorname{sgn}(\rho)\). This proves
idempotence. The product of the two projectors includes the sum
of the signs of all permutations, which is zero for
\(n\ge2\). Finally, their images are the invariant
and alternating tensors by their own averaging formulas.
\end{proof}

\subsection{Universal algebra and normalization of occupation states}
\label{vi:vi:subsec:fock-construccion}

Let \(T(\mathcal H)=\bigoplus_{n\ge0}\mathcal H^{\otimes n}\)
be the tensor algebra. The graded algebraic completion is the
quotient by the ideal generated by
\begin{equation}
 v\otimes w-(-1)^{p(v)p(w)}w\otimes v.
 \label{vi:vi:eq:ideal-fock}
\end{equation}
The universal property of the quotient means that every graded
linear map from the channels to an algebra satisfying these
relations extends uniquely to a homomorphism of the
quotient. Once declared, the exchange rule thus fixes the
product of occupations without resorting to target energies.

\begin{theorem}[Symmetric and exterior decomposition]
\label{vi:vi:thm:fock-graduado}
If \(d_0=\dim\mathcal H_{\bar0}\) and
\(d_1=\dim\mathcal H_{\bar1}\), the graded completion has
the decomposition
\begin{equation}
 \mathcal F(\mathcal H)=
 \mathcal F_+(\mathcal H_{\bar0})\otimes
 \mathcal F_-(\mathcal H_{\bar1}),
 \label{vi:vi:eq:fock-producto}
\end{equation}
\[
 \mathcal F_+(\mathcal H_{\bar0})=
 \bigoplus_{m\ge0}\operatorname{Sym}^m\mathcal H_{\bar0},
 \qquad
 \mathcal F_-(\mathcal H_{\bar1})=
 \bigoplus_{k=0}^{d_1}\bigwedge^k\mathcal H_{\bar1}.
\]
The sector of total number \(n\) is
\[
 \mathcal F_n(\mathcal H)=
 \bigoplus_{k=0}^{\min(n,d_1)}
 \operatorname{Sym}^{n-k}\mathcal H_{\bar0}
 \otimes\bigwedge^k\mathcal H_{\bar1}.
\]
\end{theorem}
\begin{proof}
Given ordered homogeneous bases, the relations
\eqref{vi:vi:eq:ideal-fock} allow the even factors to be placed first
and the odd ones afterward. Even factors commute; odd ones anticommute,
and the square of an odd one is zero because we work in characteristic
zero. Each monomial is written as
\[
 e_1^{m_1}\cdots e_{d_0}^{m_{d_0}}
 f_{i_1}\cdots f_{i_k},
 \qquad m_j\ge0,\quad i_1<\cdots<i_k.
\]
The realization through symmetric and alternating tensors makes
those monomials linearly independent. It provides the inverse
map and proves the decomposition. Grouping them by
\(m_1+\cdots+m_{d_0}+k=n\) gives the last statement.
\end{proof}

The Hilbert completion declares the normalized occupation
vectors \(|\boldsymbol m;I\rangle\) orthonormal, where
\(\boldsymbol m\in\mathbb N_0^{d_0}\) and
\(I\subset\{1,\ldots,d_1\}\). In the bosonic sector of degree
\(n=\sum_jm_j\), the tensor normalization is
\[
 |\boldsymbol m\rangle=
 \sqrt{\frac{n!}{\prod_jm_j!}}\,
 S_n(e_1^{\otimes m_1}\otimes\cdots\otimes
 e_{d_0}^{\otimes m_{d_0}}).
\]
For \(I=\{i_1<\cdots<i_k\}\), the exterior vector is
\[
 |I\rangle=\frac1{\sqrt{k!}}\sum_{\pi\in\mathfrak S_k}
 \operatorname{sgn}(\pi)
 f_{i_{\pi(1)}}\otimes\cdots\otimes f_{i_{\pi(k)}}.
\]
In the first formula there are \(n!/\prod m_j!\) distinct tensors;
in the second there are \(k!\). Counting them proves that both norms are one.
Degree zero is \(\mathbb C|0\rangle\); this occupation vacuum
is not identified with the constitutive chart of the electromagnetic vacuum
or with a vacancy in the TPK record.

\subsection{Fermionic creation and derivation of the CAR relations}
\label{vi:vi:subsec:car-construccion}

Let \(I\) be an ordered subset of odd modes and
\(r_j(I)=\#\{i\in I:i<j\}\). We define
\begin{align}
 a_j^\dagger|I\rangle&=
 \begin{cases}
 0,&j\in I,\\
 (-1)^{r_j(I)}|I\cup\{j\}\rangle,&j\notin I,
 \end{cases}\notag\\
 a_j|I\rangle&=
 \begin{cases}
 (-1)^{r_j(I)}|I\setminus\{j\}\rangle,&j\in I,\\
 0,&j\notin I.
 \end{cases}
 \label{vi:vi:eq:creacion-fermi}
\end{align}
The sign counts the transpositions needed to insert or remove
the mode while maintaining the order of the exterior basis. It does not use its mass,
its charge or the sign of a published scalar value. The matrices of
both maps are adjoints and their norms are \(1\), except for the
trivial case without odd modes.

\begin{theorem}[Canonical anticommutation relations]
\label{vi:vi:thm:car-construidas}
The operators in \eqref{vi:vi:eq:creacion-fermi} satisfy
\[
 \{a_i,a_j^\dagger\}=\delta_{ij}I,\qquad
 \{a_i,a_j\}=0,\qquad
 \{a_i^\dagger,a_j^\dagger\}=0.
\]
These are exact relations on the entirety of
\(\mathcal F_-(\mathcal H_{\bar1})\).
\end{theorem}
\begin{proof}
For \(i=j\), the product \(a_j^\dagger a_j\) is the identity
on vectors containing \(j\) and zero on the remaining ones;
\(a_ja_j^\dagger\) has the complementary actions. Their sum
is \(I\). Two insertions or two removals of the same mode give zero.
For \(i\ne j\), if any operation is blocked, both terms
of the corresponding anticommutator are zero. If both compositions
are admissible, they produce the same set of occupations. Changing
the order changes by one the number of modes preceding
one of the insertions or removals; the signs are opposite. This
proves the three anticommutators on each basis vector.
\end{proof}

\begin{corollary}[Exclusion and occupation of a complete mode]
\label{vi:vi:cor:pauli-ocupacion}
For \(N_j=a_j^\dagger a_j\),
\[
 N_j^*=N_j,\qquad N_j^2=N_j,\qquad
 \operatorname{Spec}(N_j)\subset\{0,1\}.
\]
The same complete mode admits at most one fermionic occupation.
\end{corollary}
\begin{proof}
The CAR give
\[
 N_j^2=a_j^\dagger a_ja_j^\dagger a_j
 =a_j^\dagger(I-a_j^\dagger a_j)a_j=N_j,
\]
because \(a_j^2=(a_j^\dagger)^2=0\). Self-adjointness is immediate;
an eigenvalue of a self-adjoint idempotent equals \(0\) or \(1\).
\end{proof}

The proof constructs CAR before deducing Pauli: it does not introduce
exclusion as an additional restriction on a particle table.
It also explains the real generators used in
proposition~\ref{vi:vi:prop:car-real}: that proposition now acts
on the explicit exterior algebra. Its Majorana interpretation
preserves the conditions of section~\ref{vi:vi:sec:topologia}.

\subsection{Bosonic creation, CCR and finite-particle domain}
\label{vi:vi:subsec:ccr-construccion}

On the normalized bosonic basis we define
\begin{equation}
 b_j^\dagger|\boldsymbol m\rangle=
 \sqrt{m_j+1}\,|\boldsymbol m+\boldsymbol e_j\rangle,
 \qquad
 b_j|\boldsymbol m\rangle=
 \sqrt{m_j}\,|\boldsymbol m-\boldsymbol e_j\rangle,
 \label{vi:vi:eq:creacion-bose}
\end{equation}
with value zero in the second expression when \(m_j=0\).
The factors come from the normalization of symmetric tensors.
The common domain \(\mathcal D_{\mathrm{fin}}\) is the algebraic direct
sum of the number sectors: vectors have only a
finite number of nonzero occupation components. It is dense and
invariant under the preceding maps.

\begin{theorem}[Canonical commutation relations]
\label{vi:vi:thm:ccr-construidas}
On \(\mathcal D_{\mathrm{fin}}\),
\[
 [b_i,b_j^\dagger]=\delta_{ij}I,\qquad
 [b_i,b_j]=[b_i^\dagger,b_j^\dagger]=0.
\]
The number operators \(N_j^B=b_j^\dagger b_j\) satisfy
\(N_j^B|\boldsymbol m\rangle=m_j|\boldsymbol m\rangle\).
\end{theorem}
\begin{proof}
For the same mode,
\(b_jb_j^\dagger|\boldsymbol m\rangle=(m_j+1)|\boldsymbol m\rangle\)
and \(b_j^\dagger b_j|\boldsymbol m\rangle=m_j|\boldsymbol m\rangle\).
The difference is the identity. For distinct modes, the factors
affect different occupation coordinates, so they commute.
The number formulas follow from consecutively applying the two
operations of \eqref{vi:vi:eq:creacion-bose}.
\end{proof}

If there is any even mode, \(b_j\) and \(b_j^\dagger\) are unbounded
operators in the Hilbert completion: their norms on
\(|m_j\rangle\) grow as roots of \(m_j\). The preceding
equalities therefore have an explicit domain and are not presented
as identities between finite matrices on the entire bosonic space.
Fermionic operators act on the second factor of
\eqref{vi:vi:eq:fock-producto}; bosonic ones, on the first. On
the common domain, the operators of those two factors commute.

\begin{proposition}[Exact defect of a hard bosonic cutoff]
\label{vi:vi:prop:corte-bosonico}
In one mode, let \(P_N\) be the projection onto
\(|0\rangle,\ldots,|N\rangle\) and let \(b^{(N)}=P_NbP_N\).
On that finite space,
\[
 [b^{(N)},(b^{(N)})^\dagger]
 =I-(N+1)|N\rangle\langle N|.
\]
\end{proposition}
\begin{proof}
For \(m<N\), creation and annihilation preserve the calculation
\((m+1)-m=1\). At \(|N\rangle\), truncated creation is
zero and the reverse product equals \(N\); the commutator equals
\(-N=1-(N+1)\). These are all the basis vectors.
\end{proof}

The correction resides at the last level of the cutoff, not in the occupation
law. A finite calculation can verify the CCR away from that
boundary and must preserve the boundary term when it includes the
terminal level. This distinction is indispensable so that a
matrix certificate does not confuse a truncation with the
full bosonic completion.

\subsection{Identity of complete modes and Pauli exclusion principle}
\label{vi:vi:subsec:pauli-etiquetas-hojas}

Exclusion refers to the same vector of a complete mode. Two
channels can share a cell, mass or reduced signature and remain
linearly independent through their spin, sheet, color,
flavor or another coordinate preserved by transport. Those two
vectors admit a nonzero exterior product. If a physical
quotient identifies them, that quotient must act before asserting
how many independent modes there are.

\begin{proposition}[Independence criterion for double occupation]
\label{vi:vi:prop:pauli-gram}
For two vectors of a channel space,
\[
 \|u\wedge v\|^2
 =\|u\|^2\|v\|^2-|\langle u,v\rangle|^2.
\]
In particular, two orthonormal modes produce an exterior state
of norm one and two proportional vectors produce zero.
\end{proposition}
\begin{proof}
The degree-two exterior realization is
\((u\otimes v-v\otimes u)/\sqrt2\). Expanding its norm gives
the formula. Equality in Cauchy--Schwarz characterizes linear
dependence; orthonormality gives squared norm one.
\end{proof}

For an orbital channel \(u\) and two orthogonal spin
components \(\xi_+,\xi_-\), the modes
\(u\otimes\xi_+\) and \(u\otimes\xi_-\) are distinct and their
exterior product can be occupied. There is no double occupation
of the same mode: there is occupation of two distinct components
of the representation. The same applies to two sheets when they
remain independent components of the realization.

The central electronic word may be fixed by \(C_M\), while
the charge representation distinguishes electron and positron. That
coincidence of words does not identify their complete modes or
decide their occupation. The word involution, cellular
inversion and exchange of Fock factors preserve the
differentiated domains of section~\ref{vi:vi:sec:conjugacion}.

If \(r\) orthonormal fermionic modes have the same value
of a partial observable, their total occupation can take values
\(0,1,\ldots,r\). Each \(N_j\) remains a projector;
the sum \(\sum_{j=1}^rN_j\) need not be one. This
consequence distinguishes spectral degeneracy and exclusion of
a mode, two properties that an isolated mass figure does not separate.

\subsection{Finite dimensions and transport by refinement}
\label{vi:vi:subsec:fock-refinamiento}

\begin{proposition}[Census of occupations]
\label{vi:vi:prop:censo-fock}
For \(d\) modes,
\[
 \dim\bigwedge^n\mathbb C^d=\binom dn,
 \qquad
 \dim\operatorname{Sym}^n\mathbb C^d=\binom{d+n-1}{n}
 \quad(d\ge1).
\]
The first dimension is zero for \(n>d\). In the graded
space, the dimensions of total number \(n\) are obtained
by summing the products of both dimensions for
\(k=0,\ldots,\min(n,d_1)\).
\end{proposition}
\begin{proof}
An exterior basis is chosen through subsets of \(n\) distinct
modes. A symmetric basis is chosen through
\((m_1,\ldots,m_d)\in\mathbb N_0^d\) with sum \(n\).
Placing \(d-1\) separators among \(n\) occupations
counts those weak compositions and gives the second binomial coefficient.
The graded formula follows from \eqref{vi:vi:eq:fock-producto} through
direct sums and tensor products.
\end{proof}

Let \(T:\mathcal H_K\to\mathcal H_L\) be the declared refinement
intertwiners, preserving the grading and the
corresponding channel data. Each induces
\begin{equation}
 \Gamma(T)=\bigoplus_{n\ge0}T^{\otimes n}|_{\mathcal F_n},
 \qquad T^{\otimes0}=I_{\mathbb C}.
 \label{vi:vi:eq:segunda-cuantizacion-transporte}
\end{equation}

\begin{theorem}[Compatibility of completion with transport]
\label{vi:vi:thm:fock-naturalidad}
If \(T\) preserves degree, then
\(T^{\otimes n}P_n^p=P_n^{p'}T^{\otimes n}\).
If also \(\|T\|\le1\), \(\Gamma(T)\) is a bounded
operator of norm at most \(1\) on Fock and
\[
 \Gamma(T_2T_1)=\Gamma(T_2)\Gamma(T_1),\qquad
 \Gamma(I)=I.
\]
Consequently, compatible refinement squares
of channels induce compatible squares of occupations.
\end{theorem}
\begin{proof}
Degree preservation makes \(T^{\otimes n}\) commute
with each adjacent exchange; by
proposition~\ref{vi:vi:prop:permutaciones-graduadas}, it commutes
with the permutations and their average. The restrictions
are thus well defined. Its norm is bounded by
\(\|T\|^n\le1\), so the direct sum is bounded.
Identities and compositions are verified at each degree
and, by continuity, in the completion.
\end{proof}

Contractivity is a sufficient condition that allows the bosonic
sum to be treated without additional restrictions. At finite degrees,
the powers exist for any linear map, but their
sum is not always bounded: for one bosonic mode and \(T=2I\),
\(\Gamma(T)|n\rangle=2^n|n\rangle\). In the finite exterior
sector there are only \(d_1+1\) degrees, so any
linear map of that sector has a bounded finite second
quantization. The difference is one of domain, not causality.

\begin{proposition}[Transport of creations and annihilations]
\label{vi:vi:prop:fock-transporte-adjunto}
Let \(T:\mathcal H_K\to\mathcal H_L\) be a bounded transport
preserving degree. For \(f\in\mathcal H_K\) and
\(g\in\mathcal H_L\), on the algebraic finite-particle
domain the following hold
\begin{align}
 \Gamma(T)c_K^\dagger(f)
 &=c_L^\dagger(Tf)\Gamma(T),
 \label{vi:vi:eq:fock-transporte-creacion}\\
 c_L(g)\Gamma(T)
 &=\Gamma(T)c_K(T^*g).
 \label{vi:vi:eq:fock-transporte-aniquilacion}
\end{align}
The subscripts indicate the fiber on which each operator acts.
If \(T\) is isometric, the second identity specializes to
\[
 c_L(Tf)\Gamma(T)=\Gamma(T)c_K(f),
 \qquad
 N_{L,Tf}\Gamma(T)=\Gamma(T)N_{K,f},
\]
where \(N_{K,f}=c_K^\dagger(f)c_K(f)\) and, for a unit mode,
\(N_{L,Tf}\) is its transported occupation operator.
\end{proposition}
\begin{proof}
In a symmetric or exterior product, applying \(T\) to each factor
commutes with inserting \(f\), which becomes \(Tf\). This
proves \eqref{vi:vi:eq:fock-transporte-creacion}. Annihilation
contracts each factor with the observed mode. For each factor \(v\),
\[
 \langle g,Tv\rangle=\langle T^*g,v\rangle.
\]
The exterior, or graded, signs are preserved because \(T\)
preserves degree; the symmetric normalizations are equal on
both sides. The sum of contractions proves
\eqref{vi:vi:eq:fock-transporte-aniquilacion}. If \(T^*T=I\), we
substitute \(g=Tf\). Composing the creation
and annihilation identities then gives the occupation identity. All operations
considered preserve the finite-particle domain; when
\(\|T\|\le1\), \(\Gamma(T)\) also has the bounded extension
of theorem~\ref{vi:vi:thm:fock-naturalidad}.
\end{proof}

The general identity preserves the adjoint \(T^*\) and is valid without
isometry. The need for the isometric condition appears when
transporting the same mode. For example, if \(T=tI\), with
\(0<t<1\), for a unit mode \(f\) and the vacuum \(\Omega\),
\[
 c(tf)\Gamma(T)f=t^2\Omega,
 \qquad \Gamma(T)c(f)f=\Omega.
\]
The general formula correctly gives
\(\Gamma(T)c(T^*tf)f=t^2\Omega\). This preserves
transport by contractions, with its explicit adjoint, and the
isometric specialization of occupation.

\subsection{Exchange symmetry and record of a composite}
\label{vi:vi:subsec:pauli-composicion}

The composition of section~\ref{vi:vi:subsec:composicion-registro}
acts first on histories and their boundaries. Fock organizes
channel occupation; \(\mathfrak C_{12}\) organizes the
link record. These are coordinated operations that do not
replace each other. In particular, antisymmetrization does not
suppress memory of the coupling order.

Let \(\mathcal D\subset\mathcal H^{\otimes n}\) be the space
generated by constituent expressions admissible for an
interface. The linearization of composition is written
\(C_{\mathcal T}:\mathcal D\to\mathcal H_{\mathrm{comp}}\).
When a finite group \(G\) of exchanges preserves
admissibility and the boundary, the concrete compatibility is
\[
 C_{\mathcal T}U(g)=V(g)C_{\mathcal T},\qquad g\in G,
\]
where \(U,V\) are unitary representations and \(V\) transports
the composite records, including their sheet and memory data.
The equation is the exchange
condition of that realization; it is not inferred from the mere
census of expressions.

\begin{proposition}[Descent of composition to the statistical sector]
\label{vi:vi:prop:composicion-estadistica}
Let \(\chi:G\to\{+1,-1\}\) be the character of the exchange
sector. Under the preceding compatibility,
\[
 P_\chi^{\mathrm{comp}}C_{\mathcal T}
 =C_{\mathcal T}P_\chi^{\mathrm{in}},
\]
where
\[
 P_\chi^{\mathrm{in}}=\frac1{|G|}\sum_g\chi(g)U(g),
 \qquad
 P_\chi^{\mathrm{comp}}=\frac1{|G|}\sum_g\chi(g)V(g).
\]
\end{proposition}
\begin{proof}
Multiplying each intertwining equality by
\(\chi(g)/|G|\) and summing gives exactly the identity.
The averages are projectors by the same calculation as in
proposition~\ref{vi:vi:prop:proyectores-fock}.
\end{proof}

The result justifies composing within a symmetric
or alternating sector when the boundary has the required transport.
If a permutation changes the physical binding channel, the
distinct trees and their interrelations are preserved; they are not
identified through a symmetrization that forgets that boundary.

\subsubsection{Chromatic antisymmetry and remainder of the baryonic state}

In the baryonic realization,
proposition~\ref{vi:vi:prop:singlete-barion} constructs the chromatic
vector \(|\varepsilon\rangle\), which changes sign
under a transposition. After regrouping the factors of the complete
state, a factorized vector has the form
\[
 \Psi=|\varepsilon\rangle_{\mathrm{col}}
 \otimes\Phi_{\text{spin,flavor,orbital,sheet}}.
\]
The exchange of two constituents acts simultaneously
on both factors. If \(U_{\mathrm{rest}}(\pi)\Phi=\Phi\),
then the complete state changes sign and belongs to the
fermionic sector. Conversely, for a factorized vector
with the nonzero chromatic singlet, total alternation requires
the remainder to be symmetric under those same permutations.
The proof consists in factoring
\(U_{\mathrm{tot}}(\pi)=U_{\mathrm{col}}(\pi)\otimes
U_{\mathrm{rest}}(\pi)\) and canceling the chromatic sign.

The word “remainder” explicitly includes flavor and sheet.
Exchanging colors without transporting those components is not
the permutation of two complete modes. For a nonfactorizable
state, the alternating projector is applied to the complete tensor
space, without replacing it by a scalar rule.

\subsubsection{Compatibility of physical projectors}

Color, channel and statistics projectors can be combined
by a product when they commute. If \(P,Q\) are orthogonal
projectors and \([P,Q]=0\), then
\((PQ)^*=PQ\), \((PQ)^2=PQ\) and
\(\operatorname{im}(PQ)=\operatorname{im}P\cap\operatorname{im}Q\).
Inclusion in both images is direct; for a vector
in the intersection, \(PQv=v\). If they do not commute, their product
does not in general represent the projector of the two simultaneous
conditions: the realization must construct the common
subspace or the physical order of operations.

This precision preserves the meaning of the singlet
and Pauli projectors entering the composites of
section~\ref{vi:vi:sec:familias-compuestos}. Coupling
selects a sector; the composite mass is evaluated after
constructing that sector's link record.

\subsection{Action of the material operator and scope of physical identification}
\label{vi:vi:subsec:fock-alcance}

Let \(B=B^*\) be a channel operator already obtained from the
internal dynamics and preserving the grading. At degree \(n\)
it acts through
\[
 d\Gamma_n(B)=\sum_{j=1}^n
 I^{\otimes(j-1)}\otimes B\otimes I^{\otimes(n-j)}.
\]
The sum commutes with graded permutations and restricts
to \(\mathcal F_n\). If \(B\ge0\), each restriction is positive.
In an eigenbasis, its value is the sum of the occupations
times the eigenvalues of \(B\), with the occupations
allowed by the preceding proofs. This is the action
of an additive observable on independent channels;
it does not identify the operator of a bound composite with the
sum of the isolated masses of its constituents.

In the bosonic Hilbert sum, \(d\Gamma(B)\) has the
natural domain
\[
 \left\{\psi=(\psi_n):
 \sum_n\|d\Gamma_n(B)\psi_n\|^2<\infty\right\}.
\]
Each block is self-adjoint in finite dimension, and its direct
sum with that domain is self-adjoint. The binding
or open-channel operator is added only when it has been constructed
by the corresponding dynamics. Formation of a complex
pole follows the treatment of
\ref{vi:vi:subsec:resonancias-familias}.

Complete algebraic results have therefore been obtained:
symmetry projectors, normalized bases, CAR and CCR on
their domains, exclusion of a mode, occupation censuses and
compatibility with refinement intertwiners. The starting
point of these results is declared in
\eqref{vi:vi:eq:espacio-canales-fock},
\eqref{vi:vi:eq:caracter-paridad-fock} and
\eqref{vi:vi:eq:intercambio-graduado}.

When the realization of rotations satisfies
\(U(2\pi)=(-1)^p\), the same grading coordinates
restoration of orientation at \(4\pi\) and exterior
completion. That compatibility is not replaced by a
coincidence between two signs. Identification with the
relativistic spin--statistics theorem additionally requires
the Lorentzian representation, locality and positivity
of the physical realization. Braiding sectors that do not
descend to permutations remain in the codomain of
section~\ref{vi:vi:sec:topologia}; they are not forced into
this bosonic or fermionic completion.
